\begin{figure}[!htbp]
\centering
\begin{tikzpicture}[x=1cm,y=1cm]
 % Exact printed paths of original Figure 8.2; text labels translated separately.
\path[draw={rgb,1:red,0.01112383;green,0.01551843;blue,0.01744106},line width=0.46596pt] (12.844766,2.286337) ..controls (12.875814,2.718319) and (7.652225,2.758673)..(6.786787,2.746359) ..controls (5.921373,2.758673) and (0.718580,2.715446)..(0.728819,2.283476) ..controls (0.728819,2.725746) and (5.921373,2.758673)..(6.786787,2.758673) ..controls (7.652225,2.770999) and (12.886053,2.728607)..(12.844766,2.286337);
\path[draw={rgb,1:red,0.01112383;green,0.01551843;blue,0.01744106},line width=0.34799pt,fill={rgb,1:red,0.96623176;green,0.92373538;blue,0.07670710}] (8.017962,2.341030) ..controls (8.017962,3.011181) and (7.476532,3.554440)..(6.808676,3.554440) ..controls (6.140820,3.554440) and (5.599415,3.011181)..(5.599415,2.341030) ..controls (5.599415,1.670890) and (6.140820,1.127632)..(6.808676,1.127632) ..controls (7.476532,1.127632) and (8.017962,1.670890)..(8.017962,2.341030);
\path[draw={rgb,1:red,0.01112383;green,0.01551843;blue,0.01744106},line width=0.46596pt] (12.844791,2.283159) ..controls (12.875839,1.851189) and (7.652249,1.810823)..(6.786836,1.823149) ..controls (5.921398,1.810823) and (0.718605,1.854062)..(0.728856,2.286020) ..controls (0.728856,1.843762) and (5.921398,1.810823)..(6.786836,1.810823) ..controls (7.652249,1.798510) and (12.886077,1.840901)..(12.844791,2.283159);
\path[draw={rgb,1:red,0.01112383;green,0.01551843;blue,0.01744106},line width=0.34799pt,fill={rgb,1:red,0.93083084;green,0.12043946;blue,0.14123751}] (5.182437,1.823152) ..controls (5.182437,1.886045) and (5.131501,1.937031)..(5.068656,1.937031) ..controls (5.005824,1.937031) and (4.954888,1.886045)..(4.954888,1.823152) ..controls (4.954888,1.760246) and (5.005824,1.709260)..(5.068656,1.709260) ..controls (5.131501,1.709260) and (5.182437,1.760246)..(5.182437,1.823152);
\path[draw={rgb,1:red,0.01112383;green,0.01551843;blue,0.01744106},line width=0.34799pt,fill={rgb,1:red,0.93083084;green,0.12043946;blue,0.14123751}] (5.829365,1.822900) ..controls (5.829365,1.885806) and (5.778441,1.936791)..(5.715583,1.936791) ..controls (5.652739,1.936791) and (5.601815,1.885806)..(5.601815,1.822900) ..controls (5.601815,1.760006) and (5.652739,1.709021)..(5.715583,1.709021) ..controls (5.778441,1.709021) and (5.829365,1.760006)..(5.829365,1.822900);
\path[draw={rgb,1:red,0.01112383;green,0.01551843;blue,0.01744106},line width=0.34799pt,fill={rgb,1:red,0.93083084;green,0.12043946;blue,0.14123751}] (6.922452,1.823152) ..controls (6.922452,1.886045) and (6.871528,1.937031)..(6.808696,1.937031) ..controls (6.745851,1.937031) and (6.694927,1.886045)..(6.694927,1.823152) ..controls (6.694927,1.760246) and (6.745851,1.709260)..(6.808696,1.709260) ..controls (6.871528,1.709260) and (6.922452,1.760246)..(6.922452,1.823152);
\path[draw={rgb,1:red,0.01112383;green,0.01551843;blue,0.01744106},line width=0.59158pt] (0.708159,0.942702) --(5.611092,0.942702) ..controls (5.754239,0.942702) and (5.754239,0.594705)..(5.928225,0.572349) --(7.668272,0.572349) ..controls (7.842270,0.594705) and (7.842270,0.942702)..(8.024850,0.942702) --(12.888362,0.942702);
\path[draw={rgb,1:red,0.01112383;green,0.01551843;blue,0.01744106},line width=0.34799pt,fill={rgb,1:red,0.01112383;green,0.01551843;blue,0.01744106}] (0.399056,1.823141) --(0.509830,1.887176) --(0.509830,1.759105) --(0.399056,1.823141);
\path[draw={rgb,1:red,0.01112383;green,0.01551843;blue,0.01744106},line width=0.34799pt,fill={rgb,1:red,0.01112383;green,0.01551843;blue,0.01744106}] (12.972688,0.324357) --(12.861915,0.260309) --(12.861915,0.388392) --(12.972688,0.324357);
\path[draw={rgb,1:red,0.01112383;green,0.01551843;blue,0.01744106},line width=0.69598pt,fill={rgb,1:red,0.01112383;green,0.01551843;blue,0.01744106}] (0.454446,0.323470) --(0.454446,1.822916);
\path[draw={rgb,1:red,0.01112383;green,0.01551843;blue,0.01744106},line width=0.69598pt,fill={rgb,1:red,0.01112383;green,0.01551843;blue,0.01744106}] (12.845007,0.323470) --(0.454441,0.323470);

 \node[font=\small,anchor=west] at (7.96,3.2) {恒星};
 \node[font=\small] at (4.2,2.13) {行星};
 \node[font=\small,rotate=90] at (.17,1.3) {亮度};
 \node[font=\small] at (10.4,.65) {光变曲线};
 \node[font=\small,anchor=north east] at (12.8,.16) {时间};
\end{tikzpicture}
\caption{每当行星遮挡恒星的一部分，光强就会下降。如果行星绕恒星运行，这一现象便会周期性发生。（图片由 Nikola Smolenski 提供，\url{https://en.wikipedia.org/wiki/File:Planetary_transit.svg}，许可协议 CC BY-SA 3.0。为使图示更清晰并统一风格，原书对图片作了编辑。）}
\label{fig:ch08-transit}
\end{figure}
